% Generated draft: requires mathematical and source review.
\documentclass[11pt]{article}
\usepackage[margin=25mm]{geometry}
\usepackage{fontspec}
\setmainfont{Latin Modern Roman}
\usepackage{amsmath,amssymb,mathtools,unicode-math}
\setmathfont{Latin Modern Math}
\usepackage{enumitem,needspace,xurl,hyperref}
\setlength{\parindent}{0pt}
\setlength{\parskip}{5pt}
\setlength{\emergencystretch}{3em}
\newcommand{\N}{\mathbb{N}}
\newcommand{\Z}{\mathbb{Z}}
\newcommand{\Q}{\mathbb{Q}}
\newcommand{\R}{\mathbb{R}}
\newcommand{\C}{\mathbb{C}}
\newcommand{\F}{\mathbb{F}}
\newcommand{\A}{\mathbb{A}}
\newcommand{\PP}{\mathbb{P}}
\newcommand{\E}{\mathbb{E}}
\newcommand{\eps}{\varepsilon}
\newcommand{\dd}{\,\mathrm{d}}
\begin{document}
\section*{UnsolvedMath ID: 800536}
\textbf{Source problem code:} catalogue entry
\subsection{Definitions and mathematical statement}
The supplied record does not define a \emph{$k$-blocked point set}. The following is a conditional formalization using finite planar point visibility and exactly $k$ nonempty colour classes; these conventions require confirmation against the source.

Let $k\in\Z_{\geq 1}$ and let $P\subset\R^2$ be finite. For distinct $p,q\in P$, define the open segment
\[
(p,q)=\{(1-t)p+tq:0<t<1\}.
\]
The points $p$ and $q$ are \emph{visible with respect to $P$} if $(p,q)\cap P=\varnothing$. In the interpretation used here, a $k$-blocked colouring is a surjective map $c:P\to\{1,\ldots,k\}$ satisfying
\[
\forall p,q\in P,\quad p\ne q\ \Longrightarrow\ \bigl((p,q)\cap P=\varnothing\ \Longleftrightarrow\ c(p)\ne c(q)\bigr).
\]
Thus points of different colours are visible, and every pair of distinct points of the same colour has a point of $P$ in its open connecting segment. No general-position assumption is imposed.

\textbf{Conjecture 6 (as recorded, under this interpretation).} For every $k\in\Z_{\geq 1}$, every finite $P\subset\R^2$, and every $k$-blocked colouring $c$ of $P$,
\[
\forall i\in\{1,\ldots,k\},\qquad |c^{-1}(\{i\})|\leq k.
\]
The source record labels this conjecture open; that label is supplied registry information, not an independently verified mathematical or literature assessment.
\subsection{Short English statement}
Under the stated interpretation of $k$-blocked, a finite planar point set has exactly $k$ colours, and two distinct points see each other precisely when their colours differ. The conjecture says that each colour occurs on at most $k$ points. The definition of $k$-blocked is absent from the supplied extract and requires source confirmation.
\subsection{Sources}
\begin{itemize}
\item Supplied local UnsolvedMath registry record: permanent ID $800536$, source problem code \texttt{catalogue entry}, title \emph{Visibility Graphs of Points, Segments and Polygons --- Conjecture 6}. The registry records Ghosh and Goswami as proposers, the year $2010$, and status \texttt{open}.
\item Ghosh and Goswami, \emph{Unsolved Problems in Visibility Graphs of Points, Segments and Polygons} (2010), Conjecture 6. Supplied source URL: \url{https://arxiv.org/abs/1012.5187}. The supplied background records access on 2026-10-16 and extraction by \texttt{arxiv\_latex}; these are supplied provenance statements, not claims of independent access or verification here.
\end{itemize}
\end{document}
